\ifdefined\iscombined
	\lecturetitle{Partial Fraction Decomposition}
\else
\documentclass{article}
\usepackage{xparse}
\usepackage{changepage}
\usepackage{listings}
\usepackage{amsthm}
\usepackage{comment}
\usepackage{amsmath}
\usepackage{braket}
\usepackage{amsfonts}
\usepackage{sectsty}
\usepackage{hyperref}
\usepackage{fancyhdr}
\usepackage[top=1in,bottom=1in,left=1in,right=1in]{geometry}
\usepackage{float}
\usepackage{graphicx}
\usepackage{mathtools}
\usepackage{tikz}
\usetikzlibrary{arrows}

\date{
	\vspace{-.75em}
	{\large \today}
}

\title{
	\vspace*{-1.5em}
	{\Huge Lecture 3} \\
	\vspace{-1em}
	\rule{\textwidth/4}{.01em} \\
	\vspace{-.25em}
	{\Large Partial Fraction Decomposition}
	\vspace{-.75em}
}

\author{
	{\large Matthew Farah}
}

\hypersetup{
	colorlinks=true,
	linkcolor=blue,
	filecolor=magenta,
	urlcolor=blue,
	citecolor=blue,
	anchorcolor=blue
}

\fancyhead{}
\fancyhead[L]{\today}
\fancyhead[R]{Lecture 3 - Partial Fraction Decomposition}

\setlength{\parindent}{0cm}

\newcounter{example}
\renewcommand{\theexample}{\arabic{example}}

\newenvironment{example}[1][]{
	\refstepcounter{example}
	\begin{adjustwidth}{1.5em}{}
	\noindent\textbf{\ifx&#1&Example \theexample:\else Example \theexample \ (#1):\fi}
	\ignorespaces
}{
	\vspace{.5em}
	\end{adjustwidth}
}

\newenvironment{solution}{
	\vspace{.5em}
	\par
	\begin{adjustwidth}{1.5em}{}
	\textbf{{Solution:}}
	\par
	\vspace{.5em}
}{
	\vspace{1em}
	\end{adjustwidth}
}

\newcounter{subexample}
\renewcommand{\thesubexample}{\alph{subexample}}

\newenvironment{subexample}{
	\setcounter{step}{0}
	\refstepcounter{subexample}
	\vspace{1em}\par\noindent
	\begin{adjustwidth}{1.5em}{}
	\noindent\textbf{(\thesubexample)}
	\ignorespaces
}{
	\par
	\vspace{1em}
	\end{adjustwidth}
}

\renewenvironment{proof}{
	\vspace{.5em}\par\noindent
	\begin{adjustwidth}{1.5em}{}
	\emph{Proof:}\par\vspace{1em}
	\setlength{\parindent}{0cm}
}{
	\end{adjustwidth}
}

\newtheorem{theorem}{Theorem}

\renewenvironment{theorem}[1][]{
	\refstepcounter{theorem}
	\vspace{1em}\par\noindent
	\begin{adjustwidth}{1.5em}{}
	\noindent\textbf{\ifx&#1&Theorem \thetheorem:\else Theorem \thetheorem \ (#1):\fi}
	\ignorespaces
}{
	\vspace{.5em}
	\end{adjustwidth}
}

\newtheorem{lemma}{Lemma}

\renewenvironment{lemma}[1][]{
	\refstepcounter{lemma}
	\vspace{1em}\par\noindent
	\begin{adjustwidth}{1.5em}{}
	\noindent\textbf{\ifx&#1&Lemma \thelemma:\else Lemma \thelemma \ (#1):\fi}
	\ignorespaces
}{
	\vspace{.5em}
	\end{adjustwidth}
}

\makeatother
\makeatletter

\newcounter{step}
\renewcommand{\thestep}{\Roman{step}}

\newenvironment{step}{
	\refstepcounter{step}
	\vspace{1em}\par\noindent
	\ignorespaces\noindent\textbf{(\thestep)}
	\ignorespaces
}{
	\par
}

\sectionfont{\fontsize{12}{12}\bfseries}
\subsectionfont{\fontsize{10}{12}\normalfont}
\subsubsectionfont{\fontsize{10}{12}\normalfont}

\begin{document}

\maketitle
\pagestyle{fancy}

\fi

\begin{itemize}
	\item The inverse Laplace transform takes functions expressed in the $s$-domain to the time-domain.
	\item A rational function is said to be proper if the degree of its numerator is less than or equal to the degree of its denominator.
	\item Partial fraction decomposition is a common technique used to express a transfer function whose inverse Laplace transform is unknown into the sum of fractions whose inverse Laplace transform is known.
	\begin{itemize}
		\item Partial fraction decomposition can only be performed on proper fractions.
	\end{itemize}
	\item Convolution in the time domain corresponds to multiplication in the Laplace domain and vice versa.
\end{itemize}

\begin{example}
	Perform partial fraction decomposition on $\frac{1}{s(s+2)}$.
\end{example}

\begin{solution}
	We can express that:

	\begin{align*}
		\frac{1}{s(s+2)} &= \frac{A}{s} + \frac{B}{s+2} \\
	\end{align*}

	Where $A, B \in \mathbb{R}$. Now, we can see that:

	\begin{align*}
		1 &= A(s+2) + B(s) \\
	\end{align*}

	By multiplying both sides of the equality by $s(s+2)$. Therefore:

	\begin{align*}
		1 &= As + 2A + Bs
	\end{align*}

	By equating coefficients, we have that:

	\begin{align*}
		2A &= 1 \Rightarrow A = \tfrac{1}{2} \\
		As + Bs &= 0 \Rightarrow Bs = -\tfrac{1}{2}s \Rightarrow B = -\tfrac{1}{2} \\
	\end{align*}

	Therefore:

	\begin{align*}
		\frac{1}{s(s+2)} &= \frac{1}{2s} - \frac{1}{2(s+2)}
	\end{align*}
\end{solution}

\begin{example}
	If $Y(s) = \frac{1}{s(s+2)^2(s+1)}$, find $y(t)$.
\end{example}

\begin{solution}
	To find $y(t)$, we must find the partial fraction decomposition of $Y(s)$ to compute its inverse Laplace transform. We have that:

	\begin{align*}
		Y(s) &= \frac{1}{s(s+2)^2(s+1)} = \frac{A}{s} + \frac{B}{s+2} + \frac{C}{(s+2)^2} + \frac{D}{s+1} \\
	\end{align*}

	Now, we can see that:

	\begin{align*}
		A &= sY(s) \Big|_{s=0} = \frac{1}{(s+2)^2(s+1)} \Big|_{s=0} = \frac{1}{(2)(1)} = \frac{1}{2} \\
		C &= (s+2)^2 Y(s) \Big|_{s=-2} = \frac{1}{s(s+1)} \Big|_{s=-2} = \frac{1}{(-2)(-2+1)} = \frac{1}{2} \\
		D &= (s+1)Y(s) \Big|_{s=-1} = \frac{1}{s(s+2)^2} \Big|_{s=-1} = \frac{1}{(-1)(-1+2)^2} = -1 \\
	\end{align*}

	Finding $B$, however, requires (by the product rule) that we find: %FIX: Needs fixing

	\begin{align*}
		B &= (s+2)^2 \tfrac{d}{ds} Y(s) \Big|_{s=-2} = \tfrac{d}{ds} \left( \frac{1}{s(s+1)} \right) \Big|_{s=-2} = \frac{1}{(s(s+1))^2} \Big|_{s=-2} \\
	\end{align*}

	\emph{Note to self: Come back and finish this.}
\end{solution}

\ifdefined\iscombined
	\newpage
\else
	\end{document}
\fi
